\documentclass[11pt,a4paper]{article}
\usepackage{fontspec}
\IfFontExistsTF{Times New Roman}{\setmainfont{Times New Roman}}{%
  \IfFontExistsTF{DejaVu Serif}{\setmainfont{DejaVu Serif}}{%
    \setmainfont{LibertinusSerif-Regular.otf}[
      BoldFont=LibertinusSerif-Bold.otf, ItalicFont=LibertinusSerif-Italic.otf,
      BoldItalicFont=LibertinusSerif-BoldItalic.otf]}}
\usepackage[russian]{babel}
\usepackage{amsmath,amssymb,amsthm,booktabs}
\newtheorem{mainthm}{Теорема}
\newtheorem{lem}{Лемма}
\renewcommand{\proofname}{Доказательство}
\usepackage[a4paper,margin=23mm,top=20mm,bottom=22mm]{geometry}
\usepackage{xcolor}
\definecolor{ink}{HTML}{193A4D}
\usepackage[unicode,colorlinks=true,urlcolor=ink,linkcolor=ink]{hyperref}
\setlength{\emergencystretch}{2em}
\setlength{\parindent}{1.1em}
\setlength{\parskip}{0.14em}
\raggedbottom
\hypersetup{pdftitle={Число положительных регулярных решений обратной задачи для цепочечной камерной модели с двумя выходами},pdfsubject={Положительная структурная идентифицируемость}}
\begin{document}
\begin{center}{\Large\bfseries\color{ink} Число положительных регулярных решений обратной задачи для цепочечной камерной модели с двумя выходами\par}\end{center}

\begin{center}Mikita Aharodni\\[2pt]25 ноября 2026 года\end{center}

\begin{abstract}В двунаправленной цепочке из \(n\ge4\) камер вещество вводится в первую камеру, а его количество измеряется в третьей и последней. Место выведения \(\ell\) известно. Максимальное число положительных регулярных наборов скоростей, дающих одни и те же выходные кривые, равно одному при \(\ell=3,n\) и двум при любом другом \(\ell\). Верхнюю оценку даёт вогнутость дистального усиления по спектральным весам матрицы Якоби. Точный пятикамерный пример и две операции удлинения цепочки доказывают достижимость оценки для внутренних камер. При \(n\ge5\) и \(\ell=n-1\) две реализации имеют разные константы выведения \(h\) на непустом открытом множестве данных. Если зафиксировать взаимно простые спектральные многочлены \(D,W\) и первое усиление \(p\), то совпадение \(h\) на двух регулярных ветвях возможно не более чем при одном значении второго усиления \(q\). Такой случай возникает уже при \(n=5\).\end{abstract}

{\small\noindent\textit{Ключевые слова.} структурная идентифицируемость, камерные модели, матрица Якоби, обратная спектральная задача, положительные прообразы.\par}

\section*{\color{ink} Постановка задачи}

Бортнер и соавторы \cite[теорема~5.2]{ref1} классифицировали локально идентифицируемые двунаправленные деревья с одним входом и одним выходом. Цепочки с несколькими выходами обсуждаются в \cite[§6]{ref2}. Нас интересует, сколько положительных регулярных наборов скоростей дают одни и те же наблюдения, когда фиксированы вход и два выхода. Смежный вопрос о степени идентифицируемости поставлен в \cite[вопрос~5.3]{ref3}.

Пусть \(n\ge4\), \(x'=Ax+e_1u\), \(x(0)=0\). Направленные скорости в цепи \(1\leftrightarrow\cdots\leftrightarrow n\) равны \(a_r>0\) на ребре \(r\to r+1\) и \(b_r>0\) на обратном. Единственная утечка скорости \(h>0\) расположена в известной камере \(\ell\). В матрице \(A\) имеем \(A_{r+1,r}=a_r,\ A_{r,r+1}=b_r\) и \(A_{rr}=-(a_r+b_{r-1}+h\mathbf1_{\{r=\ell\}})\), где \(a_n=b_0=0\). Остальные элементы равны нулю. Наблюдаются \(x_3,x_n\) для всех непрерывных входов.

Пусть \(F(\theta)\) состоит из коэффициентов \(D,W\) по возрастанию степеней, кроме их старших единичных коэффициентов, и усилений \(p,q\) из §1. Оба пространства имеют размерность \(2n-1\). Для фиксированных \(n,\ell\) положим
\[
M(n,\ell)=\max_y \#\{\theta>0:F(\theta)=y,\ \det DF(\theta)\ne0\}.
\]

\begin{mainthm}\label{thm:count} Для каждого \(n\ge4\)
\[
M(n,\ell)=
\begin{cases}
1,&\ell=3\ \text{или}\ \ell=n,\\
2,&\text{при всех остальных }\ell .
\end{cases}
\]

При каждом \(n,\ell\) максимум достигается положительными наборами скоростей с невырожденным якобианом. Когда максимум равен двум, оба решения сохраняются при достаточно малых общих возмущениях наблюдаемых данных.\end{mainthm}


\section*{\color{ink} 1. Симметризация и кофакторные тождества}

Используем обозначения из задачи:

\[
D(s)=\det(sI-A),\qquad W(s)=R_4(s),
\qquad p=a_1a_2,\qquad q=\prod_{r=1}^{n-1}a_r.
\]

Точные передаточные функции равны

\[
G_{31}(s)=\frac{pW(s)}{D(s)},\qquad
G_{n1}(s)=\frac q{D(s)}.
\]

Поскольку \(q>0\), вторая дробь не сокращается. Наблюдения определяют \(D,W,p,q\), включая случаи сокращения первой дроби.

Положим

\[
c_r=\sqrt{a_rb_r},\qquad \beta_r=c_r^2.
\]

Матрица \(A\) диагонально подобна симметричной матрице \(J\) с той же диагональю и внедиагональными элементами \(c_r\). Матрица \(T=-J\) положительно определена:

\[
v^TTv=
\sum_{r=1}^{n-1}
\bigl(\sqrt{a_r}v_r-\sqrt{b_r}v_{r+1}\bigr)^2
+h v_\ell^2>0
\quad(v\ne0).
\]

Для любой положительно определённой трёхдиагональной \(T\) с внедиагональными элементами \(-c_r<0\) определим

\[
w=T^{-1}e_\ell.
\]

Все координаты \(w\) положительны. Например, это следует из ряда Неймана для \(T^{-1}\): при достаточно большом \(z\) матрица \(I-T/z\) неотрицательна, имеет спектральный радиус меньше единицы, а связность цепи обеспечивает строгую положительность элементов обратной матрицы.

Положим

\[
\boxed{
a_r=c_r\frac{w_{r+1}}{w_r},\qquad
b_r=c_r\frac{w_r}{w_{r+1}},\qquad
h=\frac1{w_\ell}.}
\]

Тогда

\[
A=-\operatorname{diag}(w)\,T\,\operatorname{diag}(w)^{-1}
\]

имеет нужные переходы и единственную утечку в камере \(\ell\).

При фиксированном \(\ell\) получаем гладкое взаимно однозначное соответствие с гладким обратным отображением между физическими параметрами и такими матрицами \(T\). Поэтому ранг якобиана можно вычислять в спектральных координатах.

Обозначим через \(L_j(s)\) и \(R_j(s)\) определители ведущего блока \([1,j]\) и хвоста \([j,n]\) матрицы \(sI+T\). В частности,

\[
E=L_2,\qquad H=L_{n-1},\qquad
K=\prod_{r=3}^{n-1}\beta_r.
\]

Кофакторные формулы для \(T^{-1}\), вместе с телескопированием отношений \(w_{r+1}/w_r\), дают три тождества:

\[
\ell\ge3\quad\Longrightarrow\quad p=E(0),
\]

\[
\ell\le3\quad\Longrightarrow\quad
K=\frac{qW(0)}p,
\]

\[
q=H(0)-D(0)\frac{C_\ell(0)}{R_{\ell+1}(0)}.
\]

Здесь при \(\ell<n\) многочлен \(C_\ell\) — определитель блока \([\ell+1,n-1]\); при \(\ell=n\) последнее отношение считается нулём, так что \(q=H(0)\).

Для последнего равенства используется тождество дополнительных миноров

\[
H R_{\ell+1}-D C_\ell
=
L_{\ell-1}\prod_{r=\ell}^{n-1}\beta_r.
\]

При \(j\le\ell\) ведущий многочлен удовлетворяет \(L_{j-1}(0)=\prod_{r<j}a_r\), а при \(j\ge\ell\) чистый хвост даёт \(R_{j+1}(0)=\prod_{r=j}^{n-1}b_r\). Оба равенства получаются индукцией по длине из трёхчленной рекурсии определителя. Поэтому первая формула следует из \(L_2(0)\), вторая — из \(W(0)\), а для третьей достаточно разделить тождество дополнительных миноров в нуле на \(R_{\ell+1}(0)>0\).

\section*{\color{ink} 2. Спектральные координаты при фиксированных \(D,W\)}

Сначала предположим

\[
\gcd(D,W)=1.
\]

Случай общих корней будет включён в верхнюю границу отдельно.

Пусть

\[
0<\lambda_1<\cdots<\lambda_n
\]

— собственные значения \(T\), а \(\rho_i\) — спектральные веса в последней координате. Для неприводимой трёхдиагональной матрицы собственные значения просты, крайняя координата собственного вектора ненулевая, поэтому

\[
\rho_i>0,\qquad \sum_i\rho_i=1,
\qquad
\frac{H(s)}{D(s)}
=\sum_{i=1}^n\frac{\rho_i}{s+\lambda_i}.
\]

Пусть \(B\) — определитель блока \([4,n-1]\), с соглашением \(B=1\) при \(n=4\). Ещё одно тождество дополнительных миноров имеет вид

\[
WH-DB=KE.
\]

Обозначим

\[
P=KE.
\]

Это квадратный многочлен. Подстановка \(s=-\lambda_i\) даёт

\[
\boxed{
\rho_i=
\frac{P(-\lambda_i)}
{W(-\lambda_i)D'(-\lambda_i)},\qquad
\sum_i
\frac{P(-\lambda_i)}
{W(-\lambda_i)D'(-\lambda_i)}
=1.}
\tag{1}
\]

Нормировка задаёт аффинную плоскость в пространстве трёх коэффициентов \(P\). Условия \(\rho_i>0\) выделяют в ней открытую выпуклую область.

Пусть \(D,W\) — любые взаимно простые монические вещественные многочлены степеней \(n,n-3\) и \(D(s)=\prod_i(s+\lambda_i)\), где \(0<\lambda_1<\cdots<\lambda_n\). Положительная физическая реализация при заданном \(\ell\) существует тогда и только тогда, когда некоторый \(P\) степени не выше двух даёт по (1) положительные нормированные веса. Все такие \(P\) взаимно однозначно соответствуют реализациям. При непустом множестве решений нормировка является нетривиальным условием, поэтому это открытая часть аффинной плоскости.

Действительно, положительная дискретная мера

\[
\sum_i\rho_i\delta_{\lambda_i}
\]

однозначно восстанавливает матрицу Якоби посредством ортогонализации \(1,x,\ldots,x^{n-1}\). Обращение порядка координат и чередование знаков базиса дают нужную \(T\).

При этом для \(0\le k\le n-4\) из формулы Лагранжа следует

\[
\sum_i\rho_i W(-\lambda_i)(-\lambda_i)^k
=
\sum_i\frac{P(-\lambda_i)(-\lambda_i)^k}{D'(-\lambda_i)}
=0,
\]

поскольку степень числителя не превосходит \(n-2\). Единственность монического ортогонального многочлена степени \(n-3\) даёт характеристический многочлен крайнего блока в переменной \(x\), равный \((-1)^{n-3}W(-x)\). В переменной \(s\) это именно \(W(s)\).

Если её собственный многочлен в тождестве \(WH-DB=KE\) обозначить \(P_*\), то \(P_*\) и \(P\) совпадают в \(n\ge4\) узлах и имеют степень не выше двух. Следовательно, \(P_*=P\). В частности,

\[
[s^2]P=K>0.
\]

Эквивалентно можно использовать коэффициенты монического

\[
E(s)=s^2+es+f.
\]

Для этого положим

\[
\alpha_i(E)=
\frac{E(-\lambda_i)}
{W(-\lambda_i)D'(-\lambda_i)},\qquad
S(E)=\sum_i\alpha_i(E).
\]

Тогда

\[
\rho_i=\frac{\alpha_i}{S},\qquad K=\frac1S.
\tag{2}
\]

В физической точке \(\alpha_i>0\). Это сохраняется при малых изменениях коэффициентов \(D,W,E\), а реконструкция гладкая, поскольку используемые матрицы Грама положительно определены.

Следовательно,

\[
(\text{коэффициенты }D,\ \text{коэффициенты }W,\ e,f)
\]

— полная локальная система координат пространства скоростей. Поэтому невырожденность полного \(DF_{n,\ell}\) можно проверять по двум оставшимся наблюдениям на этой двумерной области.

\section*{\color{ink} 3. Вариационная формула для усиления}

Для \(0\le m\le n-1\) определим

\[
\Phi_m(\rho)=
\min_{\substack{\deg r\le m\\r(0)=1}}
\sum_{i=1}^n
\frac{\rho_i}{\lambda_i}\,r(\lambda_i)^2.
\tag{3}
\]

\begin{lem}\label{lem:energy}При указанной нормировке наблюдаемое усиление допускает представление

\[
\boxed{\frac q{D(0)}=\Phi_{n-\ell}(\rho).}
\tag{4}
\]\end{lem}

\subsection*{Доказательство вариационного представления}

Запишем

\[
r(x)=1-xv(x),\qquad \deg v<m,
\]

и положим \(z=v(T)e_n\). Тогда спектральное разложение даёт

\[
\sum_i\frac{\rho_i r(\lambda_i)^2}{\lambda_i}
=
e_n^TT^{-1}e_n-2e_n^Tz+z^TTz.
\]

В силу трёхдиагональности и ненулевых связей векторы \(v(T)e_n\) пробегают все векторы, поддержанные на последних \(m\) координатах.

При \(m=n-\ell\) минимум этой квадратичной функции равен

\[
e_n^TT^{-1}e_n
-
e_m^T\bigl(T_{[\ell+1,n]}\bigr)^{-1}e_m
=
\frac{H(0)}{D(0)}
-\frac{C_\ell(0)}{R_{\ell+1}(0)}.
\]

При \(m=0\) второе слагаемое отсутствует. Применяя тождество для \(q\), получаем (4).

\subsection*{Вогнутость и аналитичность}

При фиксированных \(\lambda_i\) функция \(\Phi_m\):

\par\noindent\(\bullet\)\enspace \textbf{вогнута по весам}, поскольку является минимумом линейных по \(\rho\) функций;
\par\noindent\(\bullet\)\enspace \textbf{вещественно-аналитична} при положительных весах.

Последнее можно увидеть непосредственно. Если

\[
\mu_j=\sum_i\rho_i\lambda_i^j,\qquad
b_j=\mu_j,\qquad G_{jk}=\mu_{j+k+1},
\quad 0\le j,k<m,
\]

то

\[
\Phi_m=\mu_{-1}-b^TG^{-1}b.
\tag{5}
\]

Матрица \(G\) положительно определена.

\section*{\color{ink} 4. Верхние границы числа регулярных прообразов}

\subsection*{Общая граница \(M\le2\)}

При \(\ell\ge3\) наблюдение \(p\) фиксирует \(E(0)\). В коэффициентах \(P\) это линейное условие:

\[
P(0)=p[s^2]P.
\]

При наличии решения это условие независимо от нормировки. Действительно, оно однородно и нетривиально, тогда как правая часть нормировки равна единице. Их пересечение — прямая, а положительная часть этой прямой — интервал.

При \(\ell\le2\) данные \(p,q,W\) фиксируют

\[
K=\frac{qW(0)}p.
\]

Следовательно,

\[
[s^2]P=K
\]

также выделяет прямую. Если это условие зависит от нормировки и оставляет всю плоскость, дифференциал пары \((p,q)\) на ней имеет ранг не выше единицы. Такие точки не являются регулярными.

Следовательно, веса всех регулярных решений лежат на одном аффинном интервале и удовлетворяют уравнению

\[
\Phi_{n-\ell}(\rho)=\frac q{D(0)}.
\]

У вогнутой аналитической функции на интервале не может быть трёх регулярных пересечений с горизонталью.

В самом деле, равенство её значений в трёх упорядоченных точках заставляет функцию быть постоянной между крайними. Аналитичность продолжает это постоянство на весь интервал. Тогда производная оставшегося наблюдения вдоль интервала равна нулю, и полный якобиан вырожден.

Отсюда \(M\le2\).

\subsection*{Утечка в камере \(3\)}

Здесь одновременно

\[
p=E(0),\qquad K=\frac{qW(0)}p.
\]

Поэтому нормировка (1) и оба наблюдения задают линейные условия на коэффициенты \(P\). Они имеют либо не более одной точки, либо множество положительной размерности. Во втором случае решения не регулярны.

Следовательно,

\[
M(n,3)\le1.
\]

\subsection*{Утечка в камере \(n\)}

Здесь

\[
p=E(0),\qquad q=H(0).
\]

Первое условие линейно записывается через \(P\), а \(H\) линейно зависит от \(P\) по спектральной формуле (1). Поэтому снова получаем линейную систему, имеющую не более одного изолированного регулярного решения:

\[
M(n,n)\le1.
\]

\subsection*{Данные с общими корнями \(D,W\)}

Предположим, что при каких-то данных с \(\gcd(D,W)\ne1\) имеется больше допустимого числа различных положительных регулярных прообразов.

Выберем конечное число таких точек и непересекающиеся положительные окрестности. По теореме об обратной функции все выбранные решения одновременно продолжаются на некоторую общую открытую окрестность данных.

В этой окрестности найдутся данные с

\[
\operatorname{Res}(D,W)\ne0,
\]

поскольку результант — ненулевой многочлен от независимых коэффициентов \(D,W\). Для них сохранятся все выбранные регулярные решения, что противоречит уже доказанной границе.

Тем самым верхние оценки доказаны и для данных с общими корнями \(D,W\).

\section*{\color{ink} 5. Достижимость при \(\ell=1,2\) для любого \(n\)}

Здесь можно дать одно явное спектральное семейство для всех длин.

Положим

\[
T_0=n,\qquad
t_j=j-\frac{n+1}{2},\qquad
\lambda_j=T_0+t_j,
\]

\[
a_*=\frac{n-3}{2},\qquad
b_*=\frac{n-1}{2}.
\]

Положительные собственные значения хвостового блока обозначим через \(\mu_j\). Выберем их равными

\[
\mu_j=T_0+j-\frac n2,\qquad j=2,\ldots,n-2.
\]

Они лежат во всех промежутках между соседними \(\lambda_j\), кроме двух крайних. Корни \(W(s)\) равны \(-\mu_j\).

В положительной спектральной переменной \(x\) обозначим

\[
\widehat D(x)=\prod_{j=1}^n(x-\lambda_j),\qquad
\widehat W(x)=\prod_{j=2}^{n-2}(x-\mu_j),
\]

\[
\Delta_j=\widehat W(\lambda_j)\widehat D'(\lambda_j).
\]

Рассмотрим

\[
\widehat E_u(x)
=(x-T_0)^2+u(x-T_0)-a_*b_*,
\qquad -1<u<1.
\]

Соответствующий многочлен в переменной \(s\) равен \(E_u(s)=\widehat E_u(-s)\).

У \(\Delta_j\) знак положителен на двух крайних узлах и отрицателен на всех остальных. Значения \(\widehat E_u(\lambda_j)\) имеют те же знаки. Поэтому

\[
\alpha_j(u)=\frac{\widehat E_u(\lambda_j)}{\Delta_j}>0.
\]

По симметрии

\[
\Delta_{n+1-j}=\Delta_j,\qquad
\sum_j\frac{t_j}{\Delta_j}=0.
\]

Следовательно, сумма \(S=\sum_j\alpha_j(u)\) не зависит от \(u\), как и \(K=1/S\). Нормированные веса аффинны по \(u\).

При \(u=1\) исчезают веса с индексами \(1,n-1\), а при \(u=-1\) — веса с индексами \(2,n\). В каждом конце остаётся ровно \(n-2\) узлов.

Если \(\ell=1\) или \(2\), то

\[
m=n-\ell\ge n-2.
\]

На каждом конце можно взять допустимый многочлен

\[
r_*(x)=
\prod_{\lambda\text{ в оставшемся носителе}}
(1-x/\lambda).
\]

Он имеет степень \(n-2\), равен единице в нуле и зануляется на оставшемся носителе. Поэтому

\[
\Phi_m(\rho(u))\longrightarrow0
\quad\text{при }u\to\pm1.
\]

Внутри интервала \(\Phi_m>0\), поскольку многочлен степени \(m<n\) с \(r(0)=1\) не может занулиться во всех \(n\) узлах.

Получена вогнутая аналитическая функция, положительная внутри и стремящаяся к нулю на обоих концах. Любой уровень строго между нулём и её максимумом имеет два пересечения с ненулевыми производными.

В этих двух точках совпадают \(D,W,K,q\), а значит, совпадает и

\[
p=\frac{qW(0)}K.
\]

Проверим невырожденность якобиана \(DF\).

Расширим семейство до

\[
\widehat E_{u,v}(x)
=(x-T_0)^2+u(x-T_0)+v.
\]

Тогда

\[
S=c_2+c_0v,
\]

где

\[
c_0=\sum_j\frac1{\Delta_j}
=
-\sum_{j=2}^{n-2}
\frac1{\widehat D(\mu_j)\widehat W'(\mu_j)}
<0.
\]

Равенство следует из суммы вычетов \(1/(\widehat D\widehat W)\); каждое произведение в знаменателях правой суммы положительно.

Поэтому

\[
K_v=-\frac{c_0}{S^2}\ne0.
\]

В выбранных точках \(q_u\ne0\), а \(K_u=0\). Якобиан пары \((K,q)\) по \((u,v)\) невырожден. Замена \((K,q)\leftrightarrow(p,q)\) обратима на положительном домене, поэтому невырожден весь \(DF_{n,\ell}\).

Итак,

\[
M(n,1)=M(n,2)=2
\]

для каждого \(n\ge4\).

\section*{\color{ink} 6. Достижимость единицы при \(\ell=3,n\)}

Используем то же семейство при

\[
u=0,\qquad v=-a_*b_*.
\]

Для \(\ell\ge3\)

\[
p=E(0)=T_0^2-T_0u+v.
\]

При \(\ell=3\) имеем \(q=Kp/W(0)\). Поскольку

\[
p_u=-T_0\ne0,\qquad K_u=0,\qquad K_v\ne0,
\]

якобиан пары \((p,q)\) невырожден.

При \(\ell=n\)

\[
\frac q{D(0)}=\Phi_0=\sum_j\frac{\rho_j}{\lambda_j}.
\]

В направлении \(dp=0\) берём \(dv/du=T_0\). Тогда

\[
\frac{d\alpha_j}{du}=\frac{\lambda_j}{\Delta_j},
\qquad
\frac{dS}{du}=T_0c_0,
\]

и

\[
\frac{d\Phi_0}{du}
=
\frac{c_0}{S}(1-T_0\Phi_0).
\]

В симметричной точке

\[
T_0\Phi_0
=
\sum_j\rho_j
\frac{T_0^2}{T_0^2-t_j^2}>1.
\]

Строгость обеспечивают положительные веса при \(t_j\ne0\). Поскольку \(c_0<0\), производная ненулевая. Значит, и здесь существует положительная регулярная точка.

Вместе с верхними границами это доказывает

\[
M(n,3)=M(n,n)=1.
\]

\section*{\color{ink} 7. Внутренние положения утечки. Пятикамерная модель}

Для оставшихся положений нужен способ переносить неоднозначность на произвольную длину.

Используем следующий достаточный признак существования двух регулярных прообразов. Зафиксируем \(D,W,p\) и положим

\[
E_t(s)=s^2+ts+p,\qquad t\in[a,c],
\]

причём все \(\alpha_i(t)>0\). Если для некоторого \(a<b<c\)

\[
\boxed{
\Phi_{n-\ell}(b)>
\max\{\Phi_{n-\ell}(a),\Phi_{n-\ell}(c)\},
}
\tag{6}
\]

то существуют два положительных регулярных решения с одинаковыми наблюдениями.

Действительно, нормировка \(\rho_i=\alpha_i/S\) даёт дробно-линейную инъективную параметризацию аффинного интервала весов. На этом интервале \(\Phi\) вогнута и аналитична. Скорость этой параметризации ненулевая: если вектор весов записан как \((A+tB)/(C+tD)\), то его производная равна \((BC-AD)/(C+tD)^2\), а равенство \(BC-AD=0\) противоречит инъективности. Горизонталь между большим из крайних значений и значением в \(b\) имеет два пересечения с ненулевыми производными.

Поскольку \(p=E(0)\) является отдельной координатой, ненулевая производная \(q\) вдоль уровня \(p\) означает невырожденность полного якобиана.

\subsection*{Начальный случай \((n,\ell)=(5,4)\)}

Возьмём

\[
D=s^5+17s^4+96s^3+206s^2+150s+24,
\]

\[
W=s^2+7s+8,\qquad p=2.
\]

Здесь

\[
D\bmod W=6s+24,\qquad
\operatorname{Res}(D,W)=-144\ne0.
\]

На спектральной прямой положим

\[
P_k=ks^2+\left(\frac{11k}{2}+6\right)s+2k
\]

и

\[
\begin{aligned}
H_k={}&s^4+\left(\frac k4+12\right)s^3
+\left(\frac{5k}{2}+38\right)s^2\\
&+\left(\frac{9k}{2}+36\right)s+k+6.
\end{aligned}
\]

Прямое умножение проверяет

\[
WH_k-D\left(s+2+\frac k4\right)=P_k.
\]

Последняя диагональ равна \(d_5=5-k/4\), поэтому

\[
q(k)=k+6-\frac{24}{5-k/4}.
\]

При \(k=12\) и \(k=8\) получаем следующие векторы. Сначала перечислены скорости \(a_r\), затем \(b_r\) и \(h\).

\[
\theta_A=(1,2,3,1;\ 3,1,2,2;\ 2),
\]

\[
\theta_B=
\left(
\frac{16}{29},\frac{29}{8},\frac94,\frac43;\
\frac{481}{232},\frac32,\frac89,3;\
\frac{16}{9}
\right).
\]

У обоих

\[
p=2,\qquad q=6,
\]

и совпадают указанные \(D,W\). Спектральные веса аффинны по \(k\), поэтому положительны на всём \([8,12]\).

Поскольку

\[
t=\frac{11}{2}+\frac6k,
\]

в трёх точках

\[
(a,b,c)=\left(6,\frac{61}{10},\frac{25}{4}\right)
\]

получаем

\[
\bigl(\Phi_1(a),\Phi_1(b),\Phi_1(c)\bigr)
=
\left(\frac14,\frac4{15},\frac14\right).
\]

Таким образом, строгий признак (6) выполнен.

Дополнительно точное вычисление полного якобиана даёт

\[
\det DF(\theta_A)=3456,\qquad
\det DF(\theta_B)=-4698,
\]

если коэффициенты \(D,W\) перечислять по возрастанию степеней, а затем брать \(p,q\).

\section*{\color{ink} 8. Продолжение конструкции на произвольную длину}

Пусть признак (6) уже выполнен для \((n,\ell)\), а \(m=n-\ell\).

Добавим большой положительный собственный узел \(L\) и новый корень хвоста \(L-\delta\):

\[
D^+(s)=(s+L)D(s),\qquad
W^+(s)=(s+L-\delta)W(s).
\]

Многочлен \(E_t=s^2+ts+p\) сохраняем. Выбираем

\[
L-\delta>\max_i\lambda_i.
\]

Новые ненормированные веса точно равны

\[
\alpha_i^+(t)=
\frac{\alpha_i(t)}
{(L-\lambda_i)(L-\delta-\lambda_i)}
\quad (i\le n),
\]

\[
\alpha_L^+(t)=
\frac{E_t(-L)}{-\delta W(-L)D(-L)}.
\tag{7}
\]

При достаточно большом \(L\) все они положительны на исходном отрезке \(t\). Значит, соответствующие объекты являются физическими цепями.

\subsection*{Переход A. \((n,\ell)\longmapsto(n+1,\ell+1)\)}

Берём

\[
\delta=L/2.
\]

Старые нормированные веса стремятся к исходным, а новый имеет порядок

\[
\rho_L^+=O(L^{6-2n}).
\]

Число \(m=n-\ell\) не меняется. В формуле (5) самый высокий нужный положительный момент имеет степень \(2m-1\). Вклад нового узла в него равен

\[
O(L^{6-2n+2m-1})
=
O(L^{5-2\ell})\longrightarrow0
\]

при \(\ell\ge3\).

Все остальные нужные моменты тоже сходятся. Поскольку предельная матрица Грама положительно определена,

\[
\Phi_m^+(t)\longrightarrow\Phi_m(t)
\]

в каждой из трёх фиксированных точек \(a,b,c\). Строгое неравенство (6) сохраняется при достаточно большом конечном \(L\).

\subsection*{Переход B. \((n,\ell)\longmapsto(n+1,\ell)\)}

Теперь степень остаточного многочлена увеличивается с \(m\) до \(m+1\).

Зафиксируем малое \(\varepsilon>0\) и положим

\[
\delta=\frac{L^{7-2n}}{\varepsilon}.
\]

Из (7) следует

\[
L^2\alpha_i^+(t)\longrightarrow\alpha_i(t),
\qquad
L^2\alpha_L^+(t)\longrightarrow\varepsilon.
\]

Следовательно, старые нормированные веса стремятся к \(\alpha_i/(S+\varepsilon)\), а новый — к \(\varepsilon/(S+\varepsilon)\).

Тогда для каждой фиксированной точки параметра выполняется предельное соотношение

\[
\boxed{
\Phi_{m+1}^+(t)
\longrightarrow
\frac{S(t)}{S(t)+\varepsilon}\Phi_m(t).
}
\tag{8}
\]



Для верхней оценки берём минимизатор \(r_*\) старой задачи и новый допустимый многочлен

\[
(1-x/L)r_*(x).
\]

Он зануляется в новом узле и стремится к \(r_*\) на старых узлах. Это даёт нужный \(\limsup\).

Для нижней оценки выбираем последовательность \(L\to\infty\), вдоль которой минимальные энергии стремятся к их \(\liminf\). На ней пусть \(r_L\) — минимизаторы новых задач. Их энергии ограничены уже полученной верхней оценкой. Веса старых узлов имеют положительные пределы, поэтому значения \(r_L\) во всех старых узлах ограничены.

Поскольку

\[
\deg r_L\le m+1<n,
\]

фиксированная вандермондова матрица старых узлов обеспечивает ограниченность всех коэффициентов \(r_L\).

Из энергии в новом узле следует

\[
r_L(L)=O(\sqrt L).
\]

После деления этого равенства на \(L^{m+1}\) получаем, что коэффициент при \(x^{m+1}\) стремится к нулю. Поэтому по подпоследовательности минимизаторы сходятся к многочлену степени не выше \(m\), равному единице в нуле.

Его энергия на старых узлах не меньше правой части (8), а энергия нового узла неотрицательна. Это даёт \(\liminf\) и доказывает (8).

Выберем достаточно малое фиксированное \(\varepsilon\), чтобы множители

\[
\frac{S(a)}{S(a)+\varepsilon},\quad
\frac{S(b)}{S(b)+\varepsilon},\quad
\frac{S(c)}{S(c)+\varepsilon}
\]

сохраняли строгий максимум в средней точке. Затем берём достаточно большой конечный \(L\). Так признак (6) переносится в \((n+1,\ell)\).

\subsection*{Завершение достижимости}

Из начального случая \((5,4)\):

1. Применяем переход A \(\ell-4\) раз и получаем \((\ell+1,\ell)\).
2. Применяем переход B \(n-\ell-1\) раз и получаем любую пару \((n,\ell)\) с \(4\le\ell<n\).

Все выбранные параметры конечны и строго положительны. Признак (6) даёт два положительных регулярных прообраза.

Это завершает доказательство достижимости для всех \(n,\ell\). Вместе с верхними границами получаем заявленную формулу теоремы~\ref{thm:count}.\hfill\(\square\)

\section*{\color{ink} 9. Точные сертификаты и границы применимости}

Для пяти камер и утечки в четвёртой камере точная проверка двух векторов §7 даёт определители \(3456\) и \(-4698\). Скрипт из комплекта воспроизводит также два четырёхкамерных примера, вариационное тождество на 30 рациональных моделях и строгие неравенства после переходов A, B, AB и BAB.

Прежний вектор наблюдений имел вид \(D(0),D(1),D(2),D(3),p,pW(0),q\). Здесь используются \(d_0,d_1,d_2,d_3,W(0),p,q\), где \(D(s)=s^4+\sum_{i=0}^3d_i s^i\). Определитель перехода из новых координат в прежние равен \(-12p\). Множитель \(12\) — определитель Вандермонда для узлов \(0,1,2,3\), а \(-p\) — определитель преобразования \((W(0),p)\mapsto(p,pW(0))\). Это объясняет обе величины и их знаки.

Спектральная параметризация §2 согласуется с леммами 2.1, 3.1 и теоремой 3.1 Castillo и Gordillo-Núñez \cite{ref4}. Авторы описывают семейство матриц Якоби, совместимых со спектрами полной матрицы и её главного подблока. Подсчёт в настоящей работе использует также ограничения на \(p,q\). Теорема охватывает положительные регулярные прообразы в указанной геометрии. Число комплексных и вырожденных решений требует отдельного анализа.

\section*{\color{ink} 10. Константа выведения из предпоследней камеры}
Два набора скоростей могут давать одинаковые наблюдения и при этом совпадать по отдельной координате. Для звезды с одним входом, одним выходом и без выведения идентифицируемость отдельных скоростей исследовали Clemens и соавторы \cite[теорема~3.2, §4]{ref5}. Huneke и Keen \cite[предложение~4.5]{ref6} доказали однозначность \(h\) в сильно связной модели с одним входом, когда единственный выход расположен в единственной камере выведения. Здесь выведение происходит в ненаблюдаемой камере \(n-1\).

\begin{mainthm}\label{thm:preterminal}
Пусть \(n\ge5\), единственная точка выведения известна и равна \(\ell=n-1\). При входе в камере \(1\) и наблюдениях \(x_3,x_n\) существует непустое евклидово открытое множество в пространстве коэффициентов \(F=(D,W,p,q)\), над которым имеются две положительные регулярные реализации с различными \(h\). Для этих реализаций обе выходные кривые совпадают при любом общем входе.

Более точно, зафиксируем \(D,W,p\) при \(\gcd(D,W)=1\). Если при некотором \(q\) имеются два положительных регулярных решения, то их скорости \(b_{n-1}\) различны. Среди всех значений \(q\) на этом срезе не более одного допускает два положительных регулярных решения с одинаковым \(h\). Такое исключение достигается уже при \(n=5\).
\end{mainthm}

\begin{proof}
Обозначим \(g=q/D(0)>0\). Для крайних спектральных весов \(\rho_i\) из §2 введём три момента
\[
 z=\sum_i\rho_i\lambda_i,\qquad
 v=\sum_i\rho_i\lambda_i^2,\qquad
 r=\sum_i\rho_i/\lambda_i .
\]
Это последние диагональные элементы матриц \(T,T^2,T^{-1}\) соответственно. В последней камере выведения нет, поэтому
\(z=b_{n-1}>0\) и \(v=z^2+a_{n-1}b_{n-1}>0\).
Из кофакторного тождества §1 для \(\ell=n-1\) получаем \(g=r-1/z\).

Для единичного постоянного входа стационарное состояние \(x^*=(-A)^{-1}e_1\) существует по положительной определённости \(T\). Здесь \(h\) есть константа первого порядка, а мгновенный поток выведения равен \(h x_{n-1}(t)\). Сумма стационарных балансов равна \(1-hx^*_{n-1}=0\), а баланс последней камеры равен \(a_{n-1}x^*_{n-1}-b_{n-1}x^*_n=0\). Поскольку \(x^*_n=G_{n1}(0)=g\), отсюда
\[
 g=r-\frac1z,\qquad
 h=\frac{a_{n-1}}{g\,b_{n-1}}=\frac{v-z^2}{g z^2}.
 \tag{9}
\]

Зафиксируем взаимно простые \(D,W\) и число \(p\), для которых существует положительная реализация. Согласно §§2 и 4, положительные реализации взаимно однозначно параметризуются открытым интервалом на аффинной прямой весов. Действительно, при \(\ell=n-1\ge4\) нормировку (1) дополняет линейное условие \(P(0)=p[s^2]P\), а каждый вес линейно зависит от \(P\). Все три момента аффинны вдоль этой прямой.

Если \(z\) постоянно, то \(g=r-1/z\) аффинно. Два решения на одном уровне вынуждают \(g\) быть постоянной на всём интервале. В этом случае решения лежат на непрерывном семействе с одинаковыми полными данными и оба нерегулярны. Следовательно, на любом срезе с регулярной парой \(z=b_{n-1}\) само служит координатой. В частности, его значения \(x,y\) в двух различных решениях различны. Запишем
\[
 r=\alpha z+\beta,\qquad v=\gamma z+\delta .
\]
Пара \((\gamma,\delta)\) отлична от нулевой, поскольку \(v>0\).

Пусть два решения имеют общие \(g\) и \(h\). По (9) оба числа \(z=x,y\) удовлетворяют уравнениям
\[
 \alpha z^2+(\beta-g)z-1=0,\qquad
 (1+gh)z^2-\gamma z-\delta=0 .
\]
Умножим первое уравнение на \(\delta\), вычтем второе и разделим на \(z>0\). Полученное линейное уравнение
\[
 (\delta\alpha-1-gh)z+\delta(\beta-g)+\gamma=0
\]
имеет два разных корня. Его свободный член равен нулю, откуда
\[
 \gamma=\delta(g-\beta).
 \tag{10}
\]
Обратно, пусть два различных решения имеют одно и то же \(g\) и выполнено (10). Подставим в линейное уравнение выше первое решение и его скорость \(h\). Так как \(z>0\), старший коэффициент \(\delta\alpha-1-gh\) также равен нулю. Второе квадратное уравнение тогда тождественно равно первому, умноженному на \(\delta\), и выполняется при втором корне с тем же \(h\). По (9) это и есть скорость выведения второго решения. Тем самым (10) даёт точный критерий равенства \(h\) для существующей пары.

Если \(\delta=0\), это противоречит \(\gamma\ne0\). При \(\delta\ne0\) единственно возможное значение есть \(g=\beta+\gamma/\delta\). Так как \(D(0)>0\) фиксировано, для \(q\) также возможно не более одного исключения.

Для каждого \(n\ge5\) теорема~\ref{thm:count} при \(\ell=n-1\) даёт две положительные регулярные обратные ветви над общей открытой окрестностью данных. Для этой подсемьи достаточно начальной пары \((5,4)\) из §7 и \(n-5\) применений перехода A из §8. В окрестности данных можно выбрать взаимно простые \(D,W\), поскольку их результант является ненулевым многочленом на пространстве монических коэффициентов. Зафиксируем эти \(D,W,p\). Сечение окрестности по \(q\) содержит открытый интервал. Выберем в нём любое значение, кроме возможного исключения (10). Обе ветви существуют и их значения \(h\) различны. По непрерывности это строгое различие сохраняется в открытой окрестности в полном пространстве данных. Остаётся проверить достижимость исключения.

\smallskip\noindent\textit{Точная пара с одинаковым \(h\).}
При \(n=5\) и \(\ell=4\) возьмём параметры в порядке \(a_1,\ldots,a_4,b_1,\ldots,b_4,h\):
\[
 \theta_1=(3,4,4,4,\;3,2,4,1,\;2),\qquad
 \theta_2=\left(\tfrac{44}{17},\tfrac{51}{11},\tfrac{10}{3},\tfrac{24}{5},
                 \tfrac{332}{187},\tfrac{11}{3},3,\tfrac65,2\right).
\]
Трёхчленная рекурсия даёт для обеих точек
\[
 D(s)=s^5+27s^4+230s^3+672s^2+560s+96,\qquad
 W(s)=s^2+11s+6,\qquad (p,q)=(12,192).
\]
Остаток \(D\) по модулю \(W\) равен \(-64(4s+3)\), а \(W(-3/4)=-27/16\), поэтому многочлены взаимно просты. Определители \(DF\) в порядке координат из постановки равны \(-2359296\) и \(3760128\). Следовательно, у двух регулярных решений \(h=2\), а скорости \(b_4\) равны \(1\) и \(6/5\). В примере из §7 константа \(h\) принимает значения \(2\) и \(16/9\) при одинаковых наблюдениях.
\end{proof}

\smallskip\noindent\textit{Следствие для дополнительного измерения.}
Пусть \(D,W\) взаимно просты и имеются два регулярных решения. Независимое точное измерение скорости \(b_{n-1}\) различает их. Измерение \(h\) различает решения вне возможного уровня (10). Эти выводы предполагают известное место выведения и точные наблюдения количеств вещества.

\section*{\color{ink} Вычислительная и формальная проверка}

Два скрипта с точной арифметикой проверяют коэффициенты \(D,W,p,q\) и определители якобианов для трёх пар параметров. Один вычисляет характеристические многочлены матриц, другой использует трёхчленную рекурсию и символьное дифференцирование. В Lean~4 проверены случаи \(M(4,2)=2\) и \(M(5,4)\ge2\), ряд спектральных тождеств для произвольной длины и алгебраический критерий (10) вместе с единственностью исключительного уровня. Доказательство существования регулярных ветвей для всех \(n\) дано в §§5–8. \href{https://github.com/qqhgdkktu/catenary-identifiability/tree/global-multiplicity-v2/global-multiplicity}{Исходные тексты и команды проверки}.

\begin{thebibliography}{9}
\bibitem{ref1} C. Bortner, E. Gross, N. Meshkat, A. Shiu, S. Sullivant. \textit{Identifiability of linear compartmental tree models and a general formula for input-output equations}. \href{https://arxiv.org/abs/2106.08487v2}{arXiv:2106.08487v2}, 2022. Теорема 5.2.

\bibitem{ref2} S. Ahmed, N. Crepeau, P. R. Dessauer Jr., A. Edozie, O. Garcia-Lopez, T. Grimsley, J. Lopez Garcia, V. Neri, A. Shiu. \textit{Identifiability of directed-cycle and catenary linear compartmental models}. \textit{SIAM Journal on Applied Dynamical Systems} \textbf{25} (2026), 304--350. \href{https://doi.org/10.1137/25M1728636}{DOI 10.1137/25M1728636}. §6.

\bibitem{ref3} N. Meshkat, A. Shiu. \textit{Structural identifiability of compartmental models: recent progress and future directions}. \href{https://arxiv.org/abs/2507.04496v1}{arXiv:2507.04496v1}, 2025. Вопрос 5.3.

\bibitem{ref4} K. Castillo, G. Gordillo-Núñez. \textit{Wendroff’s theorem beyond consecutive degrees and related inverse spectral problems}. \href{https://arxiv.org/abs/2602.20894v1}{arXiv:2602.20894v1}, 2026. Леммы 2.1, 3.1 и теорема 3.1.
\bibitem{ref5} K. Clemens, J. Martinez, A. Shiu, M. Thompson, B. Warren. \textit{Parameter identifiability of linear-compartmental mammillary models}. \href{https://arxiv.org/abs/2506.21889v2}{arXiv:2506.21889v2}, 2025. Теорема 3.2 и §4.
\bibitem{ref6} T. Huneke, T. Keen. \textit{On the Identifiability of Linear Compartmental Model Parameters}. Texas A\&M University, REU research report, 17 July 2025. \href{https://artsci.tamu.edu/mathematics/_files/_docs/reu/results/2025/huneke-keen-report.pdf}{Предложение~4.5}.
\end{thebibliography}
\end{document}
